% GENERATED FILE. Edit canonical lessons, then run npm run generate:books.
% canonical-source-hash: fnv1a64:e70e1fb14493e524
% canonical-lessons: RU-C01-privet, RU-C01-zdravstvuyte, RU-C01-da, RU-C01-net, RU-C01-spasibo, RU-C01-pozhaluysta, RU-C01-practice, RU-W01-false-friends-v-r, RU-W02-false-friends-s-n, RU-W03-new-shapes-b-d, RU-W04-privet-letters-p-i, RU-W05-privet-letters-e-t, RU-W05-privet-observe, RU-W05-privet-guided-copy, RU-W05-privet-delayed-copy, RU-W05-privet-dictation, RU-C01-checkpoint

\chapter{Greetings and courtesy}
\label{ch:russian-greetings}
\hlchaptermodality{1}

\begin{chapteropening}
\textbf{By the end of this chapter:} \emph{I can greet someone in Russian, respond courteously, and close the exchange naturally.}

Complete RU-C01-checkpoint, the chapter's closing checkpoint, and demonstrate the chapter promise: greet someone in Russian, respond courteously, and close the exchange naturally, then write its ten letters and two of its words from sound alone.
\end{chapteropening}

\section[\texorpdfstring{\ru{привет} (privét)}{привет (privét)}]{Privét — \textquotedblleft{}hi,\textquotedblright{} the everyday hello}

\label{lesson:RU-C01-privet}



\begin{quote}
Your first Russian word and Cyrillic letters. Several shapes are \textbf{false friends} that look Latin but sound different; meet them head-on.
\end{quote}

\subsection*{The letters in this word}
\emph{(If you already read Cyrillic, skim this.)} Russian is written \textbf{left to right}, like English. Six letters spell \emph{privét}; meet three of them now:

\begin{itemize}
\raggedright
  \item \textbf{\ru{п}} = \textquotedblleft{}p\textquotedblright{} — think Greek \textbf{pi} (Π), not Latin p.
  \item \textbf{\ru{р}} = \textbf{\textquotedblleft{}r\textquotedblright{}} — a rolled r. \textbf{FALSE FRIEND:} it looks like Latin \emph{p} but says \emph{r} (it's Greek \textbf{rho}).
  \item \textbf{\ru{в}} = \textbf{\textquotedblleft{}v\textquotedblright{}} — \textbf{FALSE FRIEND:} looks like Latin \emph{B} but says \emph{v}.
\end{itemize}

The other three, and the whole written word, arrive over this chapter. Say \textbf{privét} with the stress on the second syllable; Russian never writes its stress marks, so you learn the stress \emph{with} the word.

Two false friends in one short word — \textbf{\ru{р}}=r and \textbf{\ru{в}}=v — are the two you'll trip on most. Fix them now and half of Cyrillic stops fooling you.

\begin{cousinweb}
\emph{Privét} is built from \emph{pri-} (\textquotedblleft{}at, to, up against\textquotedblright{}) + the old root \emph{-vet} (\textquotedblleft{}to speak, to counsel\textquotedblright{}). The picture is \emph{a word brought to you} — a greeting handed over.

That root \emph{-vet} is quietly everywhere in Russian, and it hides a word you already know in English:

\begin{itemize}
\raggedright
  \item \emph{sovét} = \textquotedblleft{}a speaking-together\textquotedblright{} $\to$ \textbf{council, advice}. Borrowed into English as \textbf{Soviet} — the \emph{Soviet} Union was literally the \emph{Council} Union.
  \item \emph{otvét} = \textquotedblleft{}a speaking-back\textquotedblright{} $\to$ \textbf{an answer}.
  \item \emph{obét} = \textquotedblleft{}a vow.\textquotedblright{}
\end{itemize}

So \emph{privét} is a cousin of \textbf{Soviet} — the same \textquotedblleft{}speak\textquotedblright{} root, one a friendly hello, the other the most famous Russian word in English.
\end{cousinweb}

\begin{cousinweb}
Russian \emph{privét}, Spanish \emph{hola}, French \emph{salut}, German \emph{hallo}, and English \emph{hi} fill the casual slot. Russian draws the line sharply: use \emph{privét} with friends and the next lesson's formal greeting with everyone else.
\end{cousinweb}

\begin{culture}
\emph{Privét} is \textbf{informal} — warm, but only for people you'd address casually. To a stranger, an elder, a shopkeeper, or your boss, it can sound too familiar; there you reach for \emph{zdrávstvuyte}. Russian draws the formal/informal line more strictly than English does, and greetings are where you first feel it.
\end{culture}

\subsection*{Your turn}
Say these aloud:

\begin{itemize}
\raggedright
  \item p · r · i · v · ye · t $\to$ \textquotedblleft{}privét,\textquotedblright{} stress the \emph{vyet}
  \item which letter looks like Latin B but says \textquotedblleft{}v\textquotedblright{} (\ru{в}), and which looks like Latin p but says \textquotedblleft{}r\textquotedblright{} (\ru{р})
  \item \textquotedblleft{}privét\textquotedblright{} to a friend
\end{itemize}

\begin{tcolorbox}[breakable,colback=teal!4,colframe=teal!35!black,title={Before you move on}]
Which letter looks like Latin B, and which looks like Latin p? What does each really say? (\ru{в} says v, \ru{р} says r — the two false friends.) Is it formal or informal? (Informal — friends only.) And what famous English word shares its \emph{-vet} \textquotedblleft{}speak\textquotedblright{} root? (\emph{Soviet}, from \emph{sovét}, \textquotedblleft{}a council.\textquotedblright{})
\end{tcolorbox}
\section[\texorpdfstring{\ru{здравствуйте} (zdrávstvuyte)}{здравствуйте (zdrávstvuyte)}]{Zdrávstvuyte — \textquotedblleft{}be healthy,\textquotedblright{} the formal hello}

\label{lesson:RU-C01-zdravstvuyte}



\begin{quote}
You can say \emph{privét} to a friend. Now the word for everyone else — longer, older, and worth the effort.
\end{quote}

\subsection*{What to know first}
\begin{itemize}
\raggedright
  \item the \emph{privét} lesson — this is its formal opposite, and it reuses \ru{в}=v and \ru{р}=r from there.
\end{itemize}

\subsection*{The letters in this word}
\emph{(Skim if you read Cyrillic.)} It's a long word, and you already know two of its letters, \textbf{\ru{р}} and \textbf{\ru{в}}. Meet three more now:

\begin{itemize}
\raggedright
  \item \textbf{\ru{з}} = \textquotedblleft{}z\textquotedblright{} — looks like the digit \textbf{3}.
  \item \textbf{\ru{с}} = \textbf{\textquotedblleft{}s\textquotedblright{}} — \textbf{FALSE FRIEND:} looks like Latin \emph{C} but says \emph{s} (Greek \textbf{sigma}).
  \item \textbf{\ru{у}} = \textbf{\textquotedblleft{}oo\textquotedblright{}} (as in \emph{boot}) — \textbf{FALSE FRIEND:} looks like Latin \emph{y} but says \emph{oo}.
\end{itemize}

Its other letters arrive with this chapter's next words; the whole written word waits until every one of them is yours.

A pronunciation mercy: the first \emph{v} is \textbf{silent} here — natives say \emph{\textquotedblleft{}zdrás-tvuy-tye,\textquotedblright{}} dropping the \emph{v} in \emph{-vstv-}. Nobody articulates all those consonants.

\begin{quote}
\emph{zdrávstvuyte} (\textquotedblleft{}ZDRAST-vooy-tye\textquotedblright{})
\end{quote}

\begin{cousinweb}
\emph{Zdrávstvuyte} is a wish for \textbf{health}. Its core is \emph{zdoróvye} (\textquotedblleft{}health\textquotedblright{}) / \emph{zdoróv} (\textquotedblleft{}healthy\textquotedblright{}) — so the greeting literally means \textbf{\textquotedblleft{}be healthy!\textquotedblright{}} You are, politely, wishing wellness on the person in front of you.

\begin{itemize}
\raggedright
  \item \emph{-te} on the end is the \textbf{plural/polite ending} — the same \emph{-te} that makes a command respectful. Drop it and you get the informal singular \emph{zdrávstvuy}.
\end{itemize}

An honest note on cousins: Russian \emph{zdoróvye} (\textquotedblleft{}health\textquotedblright{}) is \textbf{not} related to English \emph{health} — they come from different roots (English \emph{health/whole/hale} is Germanic). The \emph{idea} is the shared one: many languages greet by wishing \textbf{health} (Latin \emph{salve} = \textquotedblleft{}be well,\textquotedblright{} the ancestor of Spanish \emph{hola}'s cousin \emph{salud}) or \textbf{peace} (Arabic \emph{salām}, Hebrew \emph{shalom}). Russian picked health.
\end{cousinweb}

\begin{grammarlens}[title={Grammar Lens — Russian's formal \emph{you} is plural}]
The \emph{-te} ending here is your first taste of a rule that runs through Russian: \textbf{politeness = plural.} To address one person formally, Russian uses the \emph{plural} \textquotedblleft{}you\textquotedblright{} (\emph{vy}) and plural verb endings — exactly like French \emph{vous} or the \emph{usted} idea in Spanish. So \emph{zdrávstvuy-\textbf{te}} is \textquotedblleft{}be healthy\textquotedblright{} said to a \emph{vy} — one respected person, addressed as if plural.
\end{grammarlens}

\subsection*{Your turn}
Say these aloud:

\begin{itemize}
\raggedright
  \item \textquotedblleft{}ZDRAST-vooy-tye\textquotedblright{} — let the first \ru{в} go silent
  \item which two false friends appear here (\ru{с} = s, \ru{у} = oo)
  \item \textquotedblleft{}zdrávstvuyte\textquotedblright{} to a stranger; \textquotedblleft{}privét\textquotedblright{} to a friend
\end{itemize}

\begin{tcolorbox}[breakable,colback=teal!4,colframe=teal!35!black,title={Before you move on}]
Say \emph{zdrávstvuyte}. What does it literally wish? (Health — \textquotedblleft{}be healthy.\textquotedblright{}) Formal or informal, and how do you know? (Formal — the polite plural \emph{-te}.) What do \ru{с} and \ru{у} each sound like? (s and oo — two false friends.)
\end{tcolorbox}
\section[\texorpdfstring{\ru{да} (da)}{да (da)}]{\ru{да} (da) — \textquotedblleft{}yes\textquotedblright{}}

\label{lesson:RU-C01-da}



\begin{quote}
Two greetings so far, one casual and one formal. The next two words are the shortest in the chapter and the most used in the language.
\end{quote}

\subsection*{The letters in this word}
\emph{(Skim if you read Cyrillic.)} Just two, both new and both heard in \emph{zdrávstvuyte}: \textbf{\ru{д}} (d, from Greek delta Δ) + \textbf{\ru{а}} (a, as in \emph{father}) — the one vowel that looks \emph{and} sounds like Latin a.

\begin{quote}
\textbf{\ru{да}} = \textbf{da} — short, clear, stressed. The easiest word in the language to read.
\end{quote}

\begin{cousinweb}
\textbf{\ru{да}} is an old Slavic \textbf{\ru{да}} that has always meant \textquotedblleft{}yes.\textquotedblright{} But it moonlights: Russians sprinkle \emph{\ru{да}} through speech as a soft \textbf{\textquotedblleft{}and / so / well…\textquotedblright{}} or a tag —

\begin{itemize}
\raggedright
  \item \emph{nu da} = \textquotedblleft{}well, yeah.\textquotedblright{}
  \item \emph{«\ru{да}?»} on the end of a sentence = \textquotedblleft{}…right?\textquotedblright{} (like French \emph{n'est-ce pas?}).
\end{itemize}

So the first word you can \emph{say} is also one of the most common little glue-words in the language.
\end{cousinweb}

\begin{cousinweb}
Every language improvised its own \textquotedblleft{}yes,\textquotedblright{} and you can hear where each one came from:

\begin{itemize}
\raggedright
  \item Russian \textbf{\ru{да}} — an old Slavic affirmative.
  \item Spanish \textbf{sí} — Latin \emph{sic}, \textquotedblleft{}thus, so.\textquotedblright{}
  \item French \textbf{oui} — Latin \emph{hoc ille}, \textquotedblleft{}that's it.\textquotedblright{}
  \item German \textbf{ja} — old Germanic \emph{ja}.
  \item Latin \textbf{ita} or \textbf{sic} — \textquotedblleft{}thus, so.\textquotedblright{}
\end{itemize}

Notice how many languages say \textquotedblleft{}yes\textquotedblright{} with a word that means \textbf{\textquotedblleft{}so / thus\textquotedblright{}} (\emph{sí}, \emph{sic}, \emph{ita}). Russian's \emph{\ru{да}} is its own thing — but its side-job as a \textquotedblleft{}so/and\textquotedblright{} word rhymes with that pattern.
\end{cousinweb}

\subsection*{Your turn}
Say these aloud:

\begin{itemize}
\raggedright
  \item \textquotedblleft{}da\textquotedblright{}
  \item \emph{nu da} — \textquotedblleft{}well, yeah\textquotedblright{}
  \item which letter is \ru{д} (d) and which is \ru{а} (a)
\end{itemize}

\begin{tcolorbox}[breakable,colback=teal!4,colframe=teal!35!black,title={Before you move on}]
\emph{Read it:} \textbf{\ru{да}} Two letters — name their sounds (d, a). Besides \textquotedblleft{}yes,\textquotedblright{} what else does \emph{\ru{да}} do in a sentence? (Works as a soft \textquotedblleft{}and / well… / right?\textquotedblright{})
\end{tcolorbox}
\section[\texorpdfstring{\ru{нет} (nyet)}{нет (nyet)}]{\ru{нет} (nyet) — \textquotedblleft{}no,\textquotedblright{} literally \textquotedblleft{}there-is-not\textquotedblright{}}

\label{lesson:RU-C01-net}



\begin{quote}
Say \emph{da}. One syllable, no trap in it. Its opposite has a trap.
\end{quote}

\subsection*{What to know first}
\begin{itemize}
\raggedright
  \item \emph{da} — its opposite, learned together as a pair.
\end{itemize}

\subsection*{The letters in this word}
\emph{(Skim if you read Cyrillic.)} Three letters: \textbf{\ru{н} · \ru{е} · \ru{т}}.

\begin{itemize}
\raggedright
  \item \textbf{\ru{н}} = \textbf{\textquotedblleft{}n\textquotedblright{}} — \textbf{FALSE FRIEND:} it looks exactly like Latin \emph{H} but says \emph{n}.
  \item \textbf{\ru{е}} = \textquotedblleft{}ye\textquotedblright{} (as in \emph{yes}) — looks like Latin e, but carries a \emph{y}-glide; it is the \emph{ye} of \emph{privét}.
  \item \textbf{\ru{т}} = \textquotedblleft{}t\textquotedblright{} — like Latin T.
\end{itemize}

\begin{quote}
\textbf{\ru{нет}} = \textbf{nyet} (\textquotedblleft{}nyeht\textquotedblright{}) — one syllable, stressed, with a soft \emph{y} onset.
\end{quote}

That \textbf{\ru{н}}=n is the last of the big four false friends (\ru{н}=n; the other three say \emph{v}, \emph{r} and \emph{s}). Learn those four and Cyrillic stops disguising itself as Latin.

\begin{cousinweb}
\textbf{\ru{нет}} is a \textbf{compression of a whole clause}. It comes from \textbf{\ru{не}} + \emph{(ye)st'} — \textbf{\textquotedblleft{}not is\textquotedblright{}}, i.e. \emph{\textquotedblleft{}there is not.\textquotedblright{}}

\begin{itemize}
\raggedright
  \item \textbf{\ru{не}} (\emph{nye}) — the negator, \textbf{\textquotedblleft{}not.\textquotedblright{}}
  \item \emph{yest'} — \textquotedblleft{}is / there is\textquotedblright{} (its ancestor is the same verb as English \emph{is}).
\end{itemize}

So Russian's \textquotedblleft{}no\textquotedblright{} is really \textquotedblleft{}\textbf{not-is}\textquotedblright{} — a denial of existence, worn down to a single syllable. (Its longer form \emph{nyétu}, \textquotedblleft{}there isn't any,\textquotedblright{} still shows the seam.)

And that little \textbf{\ru{не}} is one of the deepest words in the whole Indo-European family — the ancient negation \textbf{*ne}. You already own a dozen of its English children:

\begin{quote}
\textbf{\ru{не}} $\leftrightarrow$ English \textbf{no, not, none, never, nay, neither} $\leftrightarrow$ Latin \textbf{nōn} $\leftrightarrow$ Sanskrit \textbf{na} $\leftrightarrow$ French \textbf{non}.
\end{quote}

When you say \emph{\ru{нет}}, you are pronouncing a cousin of \emph{no} that has meant \textquotedblleft{}not\textquotedblright{} for six thousand years.
\end{cousinweb}

\begin{cousinweb}
Unlike \textquotedblleft{}yes\textquotedblright{} — where every language improvised — the world's \textquotedblleft{}no\textquotedblright{} is astonishingly \textbf{one word}, stretched across the whole family. Russian \textbf{\ru{нет}} carries \emph{\ru{не}}; English says \textbf{no} and \textbf{not}; Latin said \emph{nōn}, French says \emph{non}, Sanskrit said \emph{na}. All of them are the same ancient *\emph{ne}, and \emph{\ru{нет}} is Russia's share of it.
\end{cousinweb}

\subsection*{Your turn}
Say these aloud:

\begin{itemize}
\raggedright
  \item \textquotedblleft{}nyet\textquotedblright{}
  \item which letter looks like Latin H but says \textquotedblleft{}n\textquotedblright{} (\ru{н})
  \item \textquotedblleft{}\emph{da}… \ru{нет}\textquotedblright{} — yes, then no
\end{itemize}

\begin{tcolorbox}[breakable,colback=teal!4,colframe=teal!35!black,title={Before you move on}]
\emph{Read it:} \textbf{\ru{нет}} What clause is it worn down from? (\emph{\ru{не} + yest'}, \textquotedblleft{}not is.\textquotedblright{}) What does \ru{н} sound like — and what does it \emph{look} like? (Says n, looks like Latin H.) Name two English cousins of \emph{\ru{не}}. (no, not, never, none, nay…)
\end{tcolorbox}
\section[\texorpdfstring{\ru{спасибо} (spasíbo)}{спасибо (spasíbo)}]{\ru{спасибо} (spasíbo) — \textquotedblleft{}God save you,\textquotedblright{} now just \textquotedblleft{}thanks\textquotedblright{}}

\label{lesson:RU-C01-spasibo}



\begin{quote}
You can greet someone. Now thank them — and this word has a whole sentence folded inside it.
\end{quote}

\subsection*{What to know first}
\begin{itemize}
\raggedright
  \item \emph{privét} — to contrast \textbf{\ru{б}} (b) with the B-shaped letter you met there, which says \emph{v}.
\end{itemize}

\subsection*{The letters in this word}
\emph{(Skim if you read Cyrillic.)} Six letters: \textbf{\ru{с} · \ru{п} · \ru{а} · \ru{с} · \ru{и} · \ru{б} · \ru{о}}.

\begin{itemize}
\raggedright
  \item \textbf{\ru{с}} = \textquotedblleft{}s\textquotedblright{} (the false friend), \textbf{\ru{п}} = \textquotedblleft{}p\textquotedblright{}, \textbf{\ru{а}} = \textquotedblleft{}a\textquotedblright{} — all met already.
  \item \textbf{\ru{и}} = \textquotedblleft{}ee,\textquotedblright{} like a backwards Latin N; \textbf{\ru{о}} = \textquotedblleft{}o\textquotedblright{}.
  \item \textbf{\ru{б}} = \textbf{\textquotedblleft{}b\textquotedblright{}} — meet the partner that ends the confusion: \textbf{\ru{б}} is \emph{b}, while the very similar-looking letter in \emph{privét} is \emph{v}. One has an open belly (\ru{б} = b); the other stacks two bowls and says \emph{v}.
\end{itemize}

A sound note: the first \textbf{\ru{о}} is \emph{unstressed}, so it \textbf{reduces} toward \textquotedblleft{}a\textquotedblright{} — natives say \emph{\textquotedblleft{}spa-SEE-ba.\textquotedblright{}} Unstressed Russian \emph{\ru{о}} almost always softens to an \emph{a}-ish sound. Stress lands on \textbf{-\ru{си}-}.

\begin{quote}
\textbf{\ru{спасибо}} = \textbf{spasíbo} (\textquotedblleft{}spa-SEE-ba\textquotedblright{})
\end{quote}

\begin{cousinweb}
\textbf{\ru{спасибо}} is a \textbf{frozen prayer}. It is a worn-down contraction of \emph{spasí Bog} — \textbf{\textquotedblleft{}God save {[}you{]}\textquotedblright{}}:

\begin{itemize}
\raggedright
  \item \textbf{\ru{спаси}} — \textquotedblleft{}save!\textquotedblright{}, the command form of \emph{spasti} (\textquotedblleft{}to save, rescue\textquotedblright{}).
  \item \emph{Bog} — \textbf{\textquotedblleft{}God.\textquotedblright{}}
\end{itemize}

Over centuries \emph{spasí Bog} fused and shed its final consonant into today's \emph{\ru{спасибо}}. You are, etymologically, blessing whoever helped you.

Courtesy words are often fossilised blessings — Spanish \textbf{adiós} is \emph{a Dios}, English \textbf{goodbye} is \emph{God be with ye}. \emph{\ru{спасибо}} belongs to that family: a little \textquotedblleft{}God save you,\textquotedblright{} said so often the prayer wore invisible.
\end{cousinweb}

\begin{cousinweb}
\begin{itemize}
\raggedright
  \item Russian \textbf{\ru{спасибо}} — \textbf{\textquotedblleft{}God save you.\textquotedblright{}}
  \item Spanish \textbf{gracias} — Latin \emph{gratia}, \textquotedblleft{}grace.\textquotedblright{}
  \item French \textbf{merci} — Latin \emph{merces}, \textquotedblleft{}reward, mercy.\textquotedblright{}
  \item German \textbf{danke} — \textquotedblleft{}think of, be mindful.\textquotedblright{}
  \item English \textbf{thank you} — \textquotedblleft{}think\textquotedblright{}: I'll hold you in thought.
\end{itemize}

Every culture built \textquotedblleft{}thanks\textquotedblright{} from a different metaphor — grace, reward, mindful thought — and Russian built it from a \textbf{blessing}.
\end{cousinweb}

\begin{culture}
Because it began as \emph{spasí Bog}, older or very devout speakers sometimes still feel the \textquotedblleft{}God\textquotedblright{} in it; a few religious communities prefer \emph{blagodaryú} (\textquotedblleft{}I give a blessing,\textquotedblright{} a calque of Greek \emph{eucharistō}). But in everyday modern Russian, \emph{\ru{спасибо}} is simply \textbf{the} word for thanks, its prayer long dissolved.
\end{culture}

\subsection*{Your turn}
Say these aloud:

\begin{itemize}
\raggedright
  \item \textquotedblleft{}spa-SEE-ba\textquotedblright{} — reduce that first \ru{о} to an \emph{a}
  \item which letter is b (\ru{б}), and what the two-bowled letter of \emph{privét} says (v)
  \item \textquotedblleft{}\ru{спасибо}\textquotedblright{} — and recall the hidden \emph{spasí Bog}
\end{itemize}

\begin{tcolorbox}[breakable,colback=teal!4,colframe=teal!35!black,title={Before you move on}]
\emph{Read it:} \textbf{\ru{спасибо}} What two words is it worn down from, and what do they mean? (\emph{spasí Bog} — \textquotedblleft{}God save {[}you{]}.\textquotedblright{}) Which other courtesy words are fossilised blessings? (\emph{adiós} = to God; \emph{goodbye} = God be with ye.) Does \ru{б} say b or v? (b.)
\end{tcolorbox}
\section[\texorpdfstring{\ru{пожалуйста} (pozhálusta)}{пожалуйста (pozhálusta)}]{\ru{пожалуйста} (pozhálusta) — \textquotedblleft{}please,\textquotedblright{} and also \textquotedblleft{}you're welcome\textquotedblright{}}

\label{lesson:RU-C01-pozhaluysta}



\begin{quote}
Say \emph{spasíbo}. The word that answers it, and asks for things, is the longest you have met — and shorter than it looks when spoken.
\end{quote}

\subsection*{What to know first}
\begin{itemize}
\raggedright
  \item \emph{spasíbo} — this is its natural partner: \emph{spasíbo} $\to$ \emph{\ru{пожалуйста}} is Russia's \textquotedblleft{}thank you\textquotedblright{} $\to$ \textquotedblleft{}you're welcome.\textquotedblright{}
\end{itemize}

\subsection*{The letters in this word}
\emph{(Skim if you read Cyrillic.)} Everything here is familiar except one new letter: \textbf{\ru{пожалуйста}}.

\begin{itemize}
\raggedright
  \item \textbf{\ru{ж}} = \textbf{\textquotedblleft{}zh\textquotedblright{}} — the sound in \emph{mea\textbf{s}ure}. It looks like a little beetle with wings; no Latin letter carries this sound.
  \item The rest you know: \textbf{\ru{п}}(p) \textbf{\ru{о}}(o) \textbf{\ru{а}}(a) \textbf{\ru{л}}(l) \textbf{\ru{у}}(oo) \textbf{\ru{й}}(y) \textbf{\ru{с}}(s) \textbf{\ru{т}}(t).
\end{itemize}

A big pronunciation mercy: nobody says all four middle syllables. The written \emph{\ru{пожалуйста}} is spoken \textbf{\textquotedblleft{}pa-ZHAL-sta,\textquotedblright{}} swallowing the \emph{-\ru{уй}-} entirely, and the first \emph{\ru{о}} reduces to \emph{a}.

\begin{quote}
\textbf{\ru{пожалуйста}} = \textbf{pozhálusta}, said \emph{\textquotedblleft{}pa-ZHAL-sta\textquotedblright{}}
\end{quote}

\begin{cousinweb}
\textbf{\ru{пожалуйста}} is \textbf{\ru{пожалуй}} + a softening tag \textbf{-\ru{ста}}. \textbf{\ru{пожалуй}} is the command form of the old verb \emph{zhálovat'}, \textquotedblleft{}to grant, to bestow a favour\textquotedblright{} — from the root \textbf{\ru{жал}-}, the same root as \emph{zhal'} (\textquotedblleft{}a pity\textquotedblright{}) and \emph{zhalet'} (\textquotedblleft{}to feel for someone\textquotedblright{}).

So under the politeness sits \textbf{\textquotedblleft{}grant {[}me this{]}, do me the kindness.\textquotedblright{}} English \textbf{please} is the same idea worn down — short for \emph{if it please you}.

It also does a job English splits off: the \textbf{reply to thanks}. Where English needs a separate \textquotedblleft{}you're welcome,\textquotedblright{} Russian answers \emph{spasíbo} with \emph{\ru{пожалуйста}} — \textquotedblleft{}{[}it was{]} a favour freely given.\textquotedblright{}
\end{cousinweb}

\begin{cousinweb}
\begin{itemize}
\raggedright
  \item Russian \textbf{\ru{пожалуйста}} — \textquotedblleft{}grant {[}me{]} a favour.\textquotedblright{}
  \item French \textbf{s'il vous plaît} — \textquotedblleft{}if it pleases you.\textquotedblright{}
  \item Spanish \textbf{por favor} — \textquotedblleft{}for a favour.\textquotedblright{}
  \item English \textbf{please} — \textquotedblleft{}{[}if it{]} please {[}you{]}.\textquotedblright{}
  \item German \textbf{bitte} — \textquotedblleft{}I request,\textquotedblright{} and also \textquotedblleft{}you're welcome.\textquotedblright{}
\end{itemize}

German \emph{bitte} pulls the same double shift — one word for both \textquotedblleft{}please\textquotedblright{} and \textquotedblleft{}you're welcome.\textquotedblright{}
\end{cousinweb}

\begin{culture}
Because it is really a request for a kindness, \emph{\ru{пожалуйста}} also softens commands: add it to an instruction and a bark becomes a request. And as the answer to \emph{spasíbo}, it closes the ritual of thanks the way \emph{\textquotedblleft{}you're welcome\textquotedblright{}} does.
\end{culture}

\subsection*{Your turn}
Say these aloud:

\begin{itemize}
\raggedright
  \item \textquotedblleft{}pa-ZHAL-sta\textquotedblright{} — drop the middle \emph{-\ru{уй}-}
  \item the new sound \ru{ж} (the \emph{zh} in \emph{measure})
  \item the exchange — \textquotedblleft{}spasíbo!\textquotedblright{} … \textquotedblleft{}\ru{пожалуйста}!\textquotedblright{}
\end{itemize}

\begin{tcolorbox}[breakable,colback=teal!4,colframe=teal!35!black,title={Before you move on}]
\emph{Read it:} \textbf{\ru{пожалуйста}} What favour-root hides inside it? (\emph{zhálovat'}, \textquotedblleft{}to grant a favour\textquotedblright{} — cousin of \emph{zhal'}, \textquotedblleft{}pity.\textquotedblright{}) What two English phrases does this one word cover? (\textquotedblleft{}please\textquotedblright{} \textbf{and} \textquotedblleft{}you're welcome.\textquotedblright{}) What sound does \ru{ж} make? (\emph{zh}, as in \emph{measure}.)
\end{tcolorbox}
\section[Practice]{Chapter 1 — putting it together}

\label{lesson:RU-C01-practice}



\begin{quote}
Six words, all of them now in your ear and your mouth. The writing lessons that follow put their letters in your hand; first, prove you can say every one.
\end{quote}

\begin{sounds}
Four letters in these words look Latin and lie. You will write them next; for now, say each sound aloud:

\begin{itemize}
\raggedright
  \item the letter shaped like a Latin B says \textbf{v} — as in \emph{privét}.
  \item the letter shaped like a Latin P says \textbf{r} — as in \emph{privét}.
  \item the letter shaped like a Latin C says \textbf{s} — as in \emph{spasíbo}.
  \item the letter shaped like a Latin H says \textbf{n} — as in \emph{nyet}.
\end{itemize}

And the two vowel traps: the letter shaped like a Latin y says \textquotedblleft{}oo\textquotedblright{}, and the one shaped like a Latin e says \textquotedblleft{}ye.\textquotedblright{}
\end{sounds}

\subsection*{Your turn — say them cold}
Say each, then what it means:

\begin{itemize}
\raggedright
  \item \emph{privét} (hi)
  \item \emph{zdrávstvuyte} (hello, formal)
  \item \emph{spasíbo} (thank you)
  \item \emph{da} (yes)
  \item \emph{nyet} (no)
  \item \emph{pozhálusta} (please / you're welcome)
\end{itemize}

\subsection*{The exchange}
Run the whole ritual, formal then informal:

\begin{itemize}
\raggedright
  \item Meet a stranger: \emph{Zdrávstvuyte!}
  \item They thank you: \emph{Spasíbo!} $\to$ you answer \emph{Pozhálusta!}
  \item A friend waves: \emph{Privét!}
  \item Answer a question: \emph{Da!} … or \emph{Nyet!}
\end{itemize}

\begin{grammarlens}[title={What you've built — the thread through the chapter}]
Every word carried a hidden story:

\begin{itemize}
\raggedright
  \item \emph{privét} shares its \emph{-vet} \textquotedblleft{}speak\textquotedblright{} root with \textbf{Soviet} (\emph{sovét}, a council).
  \item \emph{zdrávstvuyte} wishes you \textbf{health}; \emph{spasíbo} blesses you (\emph{God save you}).
  \item \emph{nyet} is the ancient \textbf{*ne} you already own in \emph{no, not, never}.
  \item \emph{pozhálusta} asks for a \textbf{favour} — and doubles as \textquotedblleft{}you're welcome.\textquotedblright{}
\end{itemize}
\end{grammarlens}

\begin{tcolorbox}[breakable,colback=teal!4,colframe=teal!35!black,title={Before you move on}]
Without looking: which greeting is formal and which is informal? (\emph{zdrávstvuyte} / \emph{privét}.) Which word answers \emph{spasíbo}? (\emph{pozhálusta}.) Name the four false-friend letters and their real sounds. (The B-shape says v, the P-shape r, the C-shape s, the H-shape n.)
\end{tcolorbox}
\section[\texorpdfstring{\ru{в}, \ru{р} (v, r)}{в, р (v, r)}]{\ru{в} and \ru{р} — writing the two great false friends}

\label{lesson:RU-W01-false-friends-v-r}



\begin{quote}
You \emph{read} these two in \emph{\ru{привет}}. Now you'll \textbf{write} them. \textbf{\ru{в}} and \textbf{\ru{р}} are the letters that fool every beginner — they wear Latin faces (\emph{B} and \emph{P}) but say \emph{v} and \emph{r}. Forming them by hand is the fastest way to stop your eye from lying to you.
\end{quote}

\begin{culture}
Both came from \textbf{Greek}, not Latin, so the shape and the sound drifted apart:

\begin{itemize}
\raggedright
  \item \textbf{\ru{в}} is Greek \textbf{beta} — beta made the \textquotedblleft{}v\textquotedblright{} sound in later Greek, and Cyrillic kept it. It looks like Latin \emph{B}; it says \textbf{v}.
  \item \textbf{\ru{р}} is Greek \textbf{rho} — always an \textquotedblleft{}r.\textquotedblright{} It looks like Latin \emph{P}; it says a rolled \textbf{r}.
\end{itemize}
\end{culture}

\subsection*{Writing — break it apart, then write it}
\hlblockfigure{\detokenize{figures/RU-W01-false-friends-v-r-filmstrip.pdf}}{How these letters are written, one after another, stroke by stroke: \ru{в}, \ru{р}}

\textbf{\ru{в}} (\textquotedblleft{}v\textquotedblright{}) — \emph{a vertical, then two stacked bowls}:

\begin{enumerate}
\raggedright
  \item \textbf{the vertical} — a straight downstroke, top to bottom.
  \item \textbf{the upper bowl} — from the top of the vertical, bulge out right and back in.
  \item \textbf{the lower bowl} — from the middle, a second belly out to the base.
\end{enumerate}

Two bellies stacked on a spine — that's what tells \textbf{\ru{в}} apart from Latin \emph{B} (one belly) and from Russian \textbf{\ru{б}} (which you'll meet next).

\textbf{\ru{р}} (\textquotedblleft{}r\textquotedblright{}) — \emph{a vertical that drops below the line, then a top loop}:

\begin{enumerate}
\raggedright
  \item \textbf{the vertical} — a downstroke that sinks \textbf{below} the baseline (like the tail of Latin \emph{p}).
  \item \textbf{the loop} — a single closed bump at the \textbf{top}, on the right.
\end{enumerate}

So \textbf{\ru{р}} is drawn like a Latin lowercase \emph{p}, tail and all — but it is an \textbf{r}.

Their capitals, \textbf{\ru{В}} and \textbf{\ru{Р}}, are the same shapes drawn to full height, with \textbf{\ru{Р}} sitting on the line. They look even more like Latin \emph{B} and \emph{P}, and they still say \emph{v} and \emph{r}.

\subsection*{Your turn}
\begin{itemize}
\raggedright
  \item \emph{Say it:} trace \ru{в} — \textquotedblleft{}vertical, upper bowl, lower bowl\textquotedblright{} — and say its sound, \textquotedblleft{}v\textquotedblright{}
  \item \emph{Say it:} trace \ru{р} — \textquotedblleft{}vertical below the line, top loop\textquotedblright{} — and say its sound, \textquotedblleft{}r\textquotedblright{}
  \item \emph{Write it:} \emph{\ru{привет}} slowly, pausing on the \ru{в} and the \ru{р}
\end{itemize}

\begin{tcolorbox}[breakable,colback=teal!4,colframe=teal!35!black,title={Before you move on}]
Draw \textbf{\ru{в}} from its three strokes — what sound? (\textquotedblleft{}v\textquotedblright{}; a vertical + two stacked bowls.) Draw \textbf{\ru{р}} — what sound, and which Greek letter? (\textquotedblleft{}r\textquotedblright{}; Greek \textbf{rho}.) Which Latin letters do they \emph{look} like, and why is that a trap? (\emph{B} and \emph{P} — they're Greek by descent, so the shapes stayed but the sounds didn't.) Next: the other two false friends, \textbf{\ru{с}} and \textbf{\ru{н}}.
\end{tcolorbox}
\section[\texorpdfstring{\ru{с}, \ru{н} (s, n)}{с, н (s, n)}]{\ru{с} and \ru{н} — the last two false friends}

\label{lesson:RU-W02-false-friends-s-n}



\begin{quote}
Two more letters that wear a Latin mask. \textbf{\ru{с}} looks like Latin \emph{C} but says \textbf{s}; \textbf{\ru{н}} looks like Latin \emph{H} but says \textbf{n}. You've already read them — \emph{\ru{с}} opens \emph{\ru{спасибо}}, \emph{\ru{н}} opens \emph{\ru{нет}} — now write them and the four-letter false-friend set (\ru{в} \ru{р} \ru{с} \ru{н}) is complete.
\end{quote}

\subsection*{Writing — break it apart, then write it}
\hlblockfigure{\detokenize{figures/RU-W02-false-friends-s-n-filmstrip.pdf}}{How these letters are written, one after another, stroke by stroke: \ru{с}, \ru{н}}

\textbf{\ru{с}} (\textquotedblleft{}s\textquotedblright{}) — \emph{a single open curve}:

\begin{enumerate}
\raggedright
  \item \textbf{the curve} — one stroke: start top-right, sweep left and down like a \emph{C}, open on the right.
\end{enumerate}

That's the whole letter — a lone crescent. It \textbf{is} the shape of Latin \emph{C}, but it comes from \textbf{Greek sigma} (the \textquotedblleft{}s\textquotedblright{} sound). Think of it as the \textquotedblleft{}s\textquotedblright{} hiding inside the English word \emph{centre} said the soft way — \emph{c} that sounds like \emph{s}.

\textbf{\ru{н}} (\textquotedblleft{}n\textquotedblright{}) — \emph{two verticals joined by a middle bar} (a Latin \emph{H} shape):

\begin{enumerate}
\raggedright
  \item \textbf{the left vertical} — a downstroke.
  \item \textbf{the cross bar} — a short horizontal joining the two uprights \textbf{in the middle} (not slanted like a Latin \emph{N}).
  \item \textbf{the right vertical} — a second downstroke, parallel to the first.
\end{enumerate}

So \textbf{\ru{н}} is drawn exactly like a capital \emph{H} — two posts and a rung — but it is an \textbf{n}. (Russian's letter for the \emph{H}-sound doesn't exist; the \textquotedblleft{}h\textquotedblright{} of foreign words is usually spelled with \textbf{\ru{г}} or \textbf{\ru{х}}.)

The capital of \textbf{\ru{с}} is \textbf{\ru{С}}, the same crescent drawn to full height, and it still says \emph{s}.

\subsection*{Script — the four false friends, together}
\noindent
\begin{tabularx}{\linewidth}{@{}>{\raggedright\arraybackslash}X>{\raggedright\arraybackslash}X>{\raggedright\arraybackslash}X>{\raggedright\arraybackslash}X@{}}
\toprule
\textbf{letter} & \textbf{looks like} & \textbf{actually says} & \textbf{from} \\
\midrule
\textbf{\ru{в}} & Latin B & \textbf{v} & Greek beta \\
\textbf{\ru{р}} & Latin P & \textbf{r} & Greek rho \\
\textbf{\ru{с}} & Latin C & \textbf{s} & Greek sigma \\
\textbf{\ru{н}} & Latin H & \textbf{n} & — \\
\bottomrule
\end{tabularx}

Learn these four and the rest of Cyrillic stops ambushing you — most other letters either look foreign (so you can't be fooled) or behave.

\subsection*{Your turn}
\begin{itemize}
\raggedright
  \item \emph{Say it:} trace \ru{с} — \textquotedblleft{}one open curve\textquotedblright{} — sound \textquotedblleft{}s\textquotedblright{}
  \item \emph{Say it:} trace \ru{н} — \textquotedblleft{}left post, middle rung, right post\textquotedblright{} — sound \textquotedblleft{}n\textquotedblright{}
  \item \emph{Write it:} \emph{\ru{спасибо}} (find the \ru{с}) and \emph{\ru{нет}} (find the \ru{н})
\end{itemize}

\begin{tcolorbox}[breakable,colback=teal!4,colframe=teal!35!black,title={Before you move on}]
Draw \textbf{\ru{с}} — one stroke, what sound and which Latin look-alike? (\textquotedblleft{}s\textquotedblright{}; looks like \emph{C}.) Draw \textbf{\ru{н}} — three strokes, what sound and look-alike? (\textquotedblleft{}n\textquotedblright{}; looks like \emph{H}.) Name all four false friends and their sounds. (\ru{в}=v, \ru{р}=r, \ru{с}=s, \ru{н}=n.) Next: two letters with \textbf{no} Latin disguise — \ru{б} and \ru{д}.
\end{tcolorbox}
\section[\texorpdfstring{\ru{б}, \ru{д} (b, d)}{б, д (b, d)}]{\ru{б} and \ru{д} — letters that hide nothing}

\label{lesson:RU-W03-new-shapes-b-d}



\begin{quote}
A relief after the false friends: \textbf{\ru{б}} (\textquotedblleft{}b\textquotedblright{}) and \textbf{\ru{д}} (\textquotedblleft{}d\textquotedblright{}) look like \textbf{nothing in Latin}, so they can't trick you — you just have to learn to draw them. You've read both already: \emph{\ru{б}} sits in \emph{\ru{спасибо}}, \emph{\ru{д}} opens \emph{\ru{да}}.
\end{quote}

\subsection*{Script — \ru{б} (\textquotedblleft{}b\textquotedblright{}) — and how it differs from \ru{в}}
\hlblockfigure{\detokenize{figures/RU-W03-new-shapes-b-d-filmstrip.pdf}}{How these letters are written, one after another, stroke by stroke: \ru{б}, \ru{д}}

\textbf{\ru{б}} and \textbf{\ru{в}} are cousins (both from Greek \textbf{beta}, which split into a \textquotedblleft{}b\textquotedblright{} letter and a \textquotedblleft{}v\textquotedblright{} letter). They share a \textbf{vertical spine}, so keep them apart by their tops:

\begin{itemize}
\raggedright
  \item \textbf{\ru{в}} (v) = spine + \textbf{two stacked bowls} (both bellies face right).
  \item \textbf{\ru{б}} (b) = spine + \textbf{one lower belly} + a \textbf{flag/hat} curving off the \textbf{top}.
\end{itemize}

Break \textbf{\ru{б}} apart — \emph{a vertical, a lower belly, a top flag}:

\begin{enumerate}
\raggedright
  \item \textbf{the vertical} — a downstroke.
  \item \textbf{the belly} — one bowl bulging right from the \textbf{lower} half.
  \item \textbf{the top flag} — a short stroke off the top, curving \textbf{right} (the \textquotedblleft{}hat\textquotedblright{} that \ru{в} doesn't have).
\end{enumerate}

That top flag is the whole difference: \textbf{\ru{б}} has a hat and one belly; \textbf{\ru{в}} has no hat and two bellies.

\subsection*{Script — \ru{д} (\textquotedblleft{}d\textquotedblright{}) — the delta on little feet}
\textbf{\ru{д}} comes from Greek \textbf{delta} (Δ), the triangle — and you can still feel the triangle in it. Break it apart — \emph{a body, then two little feet}:

\begin{enumerate}
\raggedright
  \item \textbf{the body} — a rounded/looped shape sitting on the baseline (think a small \emph{delta} or a curled \emph{g}-like body).
  \item \textbf{the feet} — two little legs dropping \textbf{below} the baseline, one at each corner of the base.
\end{enumerate}

Those two descending feet are \textbf{\ru{д}}'s signature — no Latin letter has them. Say \textquotedblleft{}delta\textquotedblright{} as you draw it and the shape sticks.

Their capitals: \textbf{\ru{Б}} keeps the belly but swaps the curling flag for a straight bar across the top; \textbf{\ru{Д}} is the delta on its feet, drawn to full height.

\subsection*{Your turn}
\begin{itemize}
\raggedright
  \item \emph{Say it:} trace \ru{б} — \textquotedblleft{}spine, lower belly, top flag\textquotedblright{} — sound \textquotedblleft{}b\textquotedblright{} — then \ru{в} for contrast (\textquotedblleft{}two bellies, no hat\textquotedblright{})
  \item \emph{Say it:} trace \ru{д} — \textquotedblleft{}body, two feet below the line\textquotedblright{} — sound \textquotedblleft{}d\textquotedblright{}, from delta
  \item \emph{Write it:} \emph{\ru{спасибо}} (find the \ru{б}) and \emph{\ru{да}} (the \ru{д})
\end{itemize}

\begin{tcolorbox}[breakable,colback=teal!4,colframe=teal!35!black,title={Before you move on}]
What one stroke separates \textbf{\ru{б}} from \textbf{\ru{в}}? (The \textbf{top flag/hat} — \ru{б} has it, and one belly; \ru{в} has two bellies, no hat.) Draw \textbf{\ru{д}} — what sound, from which Greek letter, and what's its signature? (\textquotedblleft{}d\textquotedblright{}; from \textbf{delta}; the two little feet below the line.) You can now hand-write every letter in \emph{\ru{привет}}, \emph{\ru{спасибо}}, \emph{\ru{да}}, and \emph{\ru{нет}}.
\end{tcolorbox}
\section[\texorpdfstring{\ru{п}, \ru{и} (p, i)}{п, и (p, i)}]{\ru{п} and \ru{и} — filling in \ru{привет}}

\label{lesson:RU-W04-privet-letters-p-i}



\begin{quote}
You can already hand-write the \emph{\ru{в}} and the \emph{\ru{р}} of \emph{\ru{привет}}. Now take two more of its letters: \textbf{\ru{п}} (\textquotedblleft{}p\textquotedblright{}) and \textbf{\ru{и}} (\textquotedblleft{}ee\textquotedblright{}). One is an honest Greek shape; the other is a \textbf{quiet false friend} — it looks like a Latin \emph{N}.
\end{quote}

\subsection*{Script — \ru{п} (\textquotedblleft{}p\textquotedblright{}) — the Greek gate}
\hlblockfigure{\detokenize{figures/RU-W04-privet-letters-p-i-filmstrip.pdf}}{How these letters are written, one after another, stroke by stroke: \ru{п}, \ru{и}}

\textbf{\ru{п}} is Greek \textbf{pi} (Π), the letter mathematicians write for 3.14…. It's an honest letter: it looks like a little gateway and says \textbf{\textquotedblleft{}p.\textquotedblright{}}

Break it apart — \emph{two verticals joined by a top bar}:

\begin{enumerate}
\raggedright
  \item \textbf{the left vertical} — a downstroke.
  \item \textbf{the top bar} — a short horizontal across the top.
  \item \textbf{the right vertical} — a second downstroke, parallel to the first.
\end{enumerate}

A doorway shape: two posts and a lintel. (Don't confuse it with \textbf{\ru{н}} — that one's bar sits in the \textbf{middle}, this one's sits on \textbf{top}.)

\subsection*{Script — \ru{и} (\textquotedblleft{}ee\textquotedblright{}) — the backwards N}
\textbf{\ru{и}} says \textbf{\textquotedblleft{}ee\textquotedblright{}} (as in \emph{meet}) — a vowel. But watch out: it looks almost exactly like a \textbf{Latin N with the diagonal reversed}. That's the quiet trap — your eye reads \textquotedblleft{}N,\textquotedblright{} but it's a vowel.

Break it apart — \emph{two verticals joined by a rising diagonal}:

\begin{enumerate}
\raggedright
  \item \textbf{the left vertical} — a downstroke.
  \item \textbf{the diagonal} — from the \textbf{bottom} of the left vertical \textbf{up} to the top of the right (rising left-to-right — the mirror of Latin \emph{N}, whose diagonal \textbf{falls}).
  \item \textbf{the right vertical} — a downstroke.
\end{enumerate}

So the tell between \emph{\ru{и}} and Latin \emph{N} is which way the diagonal leans: \textbf{\ru{и}} rises, \emph{N} falls. (Later you'll meet \textbf{\ru{й}} — an \emph{\ru{и}} with a little breve on top, \textquotedblleft{}ee-y.\textquotedblright{})

Their capitals, \textbf{\ru{П}} and \textbf{\ru{И}}, are the same shapes drawn to full height. So is \textbf{\ru{Н}}, the capital of \textbf{\ru{н}}, with its bar still in the middle.

\subsection*{Your turn}
\begin{itemize}
\raggedright
  \item \emph{Say it:} trace \ru{п} — \textquotedblleft{}left post, top bar, right post\textquotedblright{} — sound \textquotedblleft{}p\textquotedblright{}
  \item \emph{Say it:} trace \ru{и} — \textquotedblleft{}left post, rising diagonal, right post\textquotedblright{} — sound \textquotedblleft{}ee\textquotedblright{}
  \item \emph{Write it:} the start of \emph{\ru{привет}} — \ru{п} · \ru{р} · \ru{и} · … — three letters down
\end{itemize}

\begin{tcolorbox}[breakable,colback=teal!4,colframe=teal!35!black,title={Before you move on}]
Draw \textbf{\ru{п}} — what sound and which Greek letter? (\textquotedblleft{}p\textquotedblright{}; Greek \textbf{pi} Π.) Draw \textbf{\ru{и}} — what sound, and how do you tell it from Latin \emph{N}? (\textquotedblleft{}ee\textquotedblright{}; its diagonal \textbf{rises} left-to-right, where \emph{N}'s falls.) Which letter has its bar on top, \ru{п} or \ru{н}? (\ru{п} — top bar; \ru{н} — middle bar.) Next: \textbf{\ru{е}} and \textbf{\ru{т}}, and you'll have all of \emph{\ru{привет}}.
\end{tcolorbox}
\section[\texorpdfstring{\ru{е}, \ru{т} (e, t)}{е, т (e, t)}]{\ru{е} and \ru{т} — the last two letters of \ru{привет}}

\label{lesson:RU-W05-privet-letters-e-t}



\begin{quote}
The finish line. \textbf{\ru{е}} and \textbf{\ru{т}} are two of Cyrillic's most \textbf{honest} letters — they look like their Latin twins \emph{and} sound close to them. Learn these and you can hand-write the whole of \emph{\ru{привет}} from memory.
\end{quote}

\subsection*{Script — \ru{е} (\textquotedblleft{}ye\textquotedblright{}) — the honest vowel}
\textbf{\ru{е}} looks like a Latin \textbf{e} and sounds close to one — with a small twist: at the start of a syllable it carries a \textbf{y}-glide, \textquotedblleft{}\textbf{ye}\textquotedblright{} (as in \emph{yes}). So \emph{\ru{привет}} ends \emph{…v-\textbf{ye}-t}. (Letters that add a \emph{y}-onset like this are called \emph{iotated} — \emph{\ru{е}, \ru{ё}, \ru{ю}, \ru{я}} are the family; you'll meet the others in Chapter 2.)

Break it apart — \emph{a curve, then a middle bar} (exactly like writing a Latin \emph{e}):

\begin{enumerate}
\raggedright
  \item \textbf{the curve} — start top-right, sweep left, curl down and round.
  \item \textbf{the middle bar} — the short horizontal that closes the \emph{e}.
\end{enumerate}

No trap here: what you see is what you get.

\subsection*{Script — \ru{т} (\textquotedblleft{}t\textquotedblright{}) — Greek tau, honest too}
\textbf{\ru{т}} is Greek \textbf{tau}, and it looks and sounds like Latin \textbf{T}: a crossbar on a stem. (In \emph{handwriting} Russians often make lowercase \emph{\ru{т}} look like an \emph{m} with a bar over it, but the printed citation form is the clean T-shape you see here.)

Break it apart — \emph{a central downstroke, then a top bar}:

\begin{enumerate}
\raggedright
  \item \textbf{the downstroke} — a vertical stem.
  \item \textbf{the top bar} — a horizontal across the top, centred on the stem.
\end{enumerate}

\subsection*{Writing — now write the whole word}
\hlblockfigure{\detokenize{figures/RU-W05-privet-letters-e-t-filmstrip.pdf}}{How these letters are written, one after another, stroke by stroke: \ru{е}, \ru{т}}

You now own every letter of \textbf{\ru{привет}}:

\begin{quote}
\textbf{\ru{п}} (p) · \textbf{\ru{р}} (r) · \textbf{\ru{и}} (ee) · \textbf{\ru{в}} (v) · \textbf{\ru{е}} (ye) · \textbf{\ru{т}} (t) = \emph{privét}
\end{quote}

Two of them are false friends (\textbf{\ru{р}}=r, \textbf{\ru{в}}=v), one is a quiet one (\textbf{\ru{и}}=ee), and three are honest (\textbf{\ru{п}, \ru{е}, \ru{т}}). Write it slowly, naming each letter's sound.

\subsection*{Your turn}
\begin{itemize}
\raggedright
  \item \emph{Say it:} trace \ru{е} — \textquotedblleft{}curve, middle bar\textquotedblright{} — sound \textquotedblleft{}ye\textquotedblright{}
  \item \emph{Say it:} trace \ru{т} — \textquotedblleft{}downstroke, top bar\textquotedblright{} — sound \textquotedblleft{}t\textquotedblright{}, Greek tau
  \item \emph{Write it:} the full word — \emph{\ru{привет}} — start to finish, from memory
\end{itemize}

\begin{tcolorbox}[breakable,colback=teal!4,colframe=teal!35!black,title={Before you move on}]
Draw \textbf{\ru{е}} — what sound, and what does \emph{iotated} mean? (\textquotedblleft{}Ye\textquotedblright{}; it carries a \emph{y}-onset.) Draw \textbf{\ru{т}} — what sound and which Greek letter? (\textquotedblleft{}t\textquotedblright{}; Greek \textbf{tau}.) Spell out the letters of \emph{\ru{привет}} with their sounds. (\ru{п}-\ru{р}-\ru{и}-\ru{в}-\ru{е}-\ru{т} = p-r-i-v-ye-t.) You can now hand-write your first Russian word end to end.
\end{tcolorbox}
\section[Practice]{\ru{привет} — watch one word move}

\label{lesson:RU-W05-privet-observe}



\begin{quote}
Look at \textbf{\ru{привет}}. Point from left to right and say the six sounds slowly: \emph{p, r, ee, v, ye, t}. Every shape is already known.
\end{quote}

\subsection*{Writing — observe and trace}
Keep \textbf{\ru{привет}} visible. Trace it once with a finger or a capped pen. Pause on \textbf{\ru{р}} and say \emph{r}; pause on \textbf{\ru{в}} and say \emph{v}. Stop after one careful trace.

\begin{tcolorbox}[breakable,colback=teal!4,colframe=teal!35!black,title={Before you move on}]
Read \textbf{\ru{привет}} once. Today was observation, not a memory test.
\end{tcolorbox}
\section[\texorpdfstring{\ru{привет} (privét)}{привет (privét)}]{\ru{привет} — one supported copy}

\label{lesson:RU-W05-privet-guided-copy}



\begin{quote}
Keep \textbf{\ru{привет}} visible. Point to its two traps: \textbf{\ru{р}} says \emph{r} and \textbf{\ru{в}} says \emph{v}. You may look back after every letter.
\end{quote}

\subsection*{Writing — guided copy}
Copy \textbf{\ru{привет}} once. Use three small chunks: \textbf{\ru{при} · \ru{ве} · \ru{т}}. After each chunk, look back and compare only those letters. Then read the finished word.

\begin{tcolorbox}[breakable,colback=teal!4,colframe=teal!35!black,title={Before you move on}]
Stop after one supported copy. Hiding the model belongs to the next lesson.
\end{tcolorbox}
\section[\texorpdfstring{\ru{привет} (privét)}{привет (privét)}]{\ru{привет} — look, hide, write, compare}

\label{lesson:RU-W05-privet-delayed-copy}



\begin{quote}
Look at \textbf{\ru{привет}} for five seconds. Touch the three chunks \textbf{\ru{при} · \ru{ве} · \ru{т}}. Do not copy while the model is visible.
\end{quote}

\subsection*{Writing — delayed copy}
\begin{itemize}
\raggedright
  \item cover \textbf{\ru{привет}}
  \item wait five seconds
  \item write it once from memory
  \item uncover the model and compare one chunk at a time
\end{itemize}

Circle and repair one differing letter if needed. Repair is part of retrieval.

\begin{tcolorbox}[breakable,colback=teal!4,colframe=teal!35!black,title={Before you move on}]
Read the repaired word once. Stop after one memory attempt.
\end{tcolorbox}
\section[\texorpdfstring{\ru{привет} (privét)}{привет (privét)}]{One heard greeting — write what you know}

\label{lesson:RU-W05-privet-dictation}



\begin{quote}
Cover the answer lower on this page. You will hear one familiar informal greeting. No new Russian and no new Cyrillic shape is hiding here.
\end{quote}

\subsection*{Writing — short dictation}
\emph{Hear:} \emph{privét} — informal \textquotedblleft{}hi\textquotedblright{}

Write the Russian word from sound alone. Do not use a visible model or romanization. Now uncover and compare: \textbf{\ru{привет}}. Check \textbf{\ru{р}} and \textbf{\ru{в}} first, then the other four letters. Repair one place if needed.

\begin{tcolorbox}[breakable,colback=teal!4,colframe=teal!35!black,title={Before you move on}]
This completes a tiny writing runway, not an exam simulation. Later work must grow from one word to messages, connected composition, and timed production.
\end{tcolorbox}
\section[Practice]{Chapter 1 checkpoint — say it, know the letters, write it}

\label{lesson:RU-C01-checkpoint}



\begin{quote}
Nothing new arrives here. This checks what the chapter built: the spoken exchange, ten letters, and two words from sound alone. First: which greeting goes to a friend, and which to a stranger? (\textbf{\ru{Привет}}; \emph{zdrávstvuyte}.)
\end{quote}

\subsection*{The exchange — formal, then informal}
With a stranger:

\begin{itemize}
\raggedright
  \item \emph{Zdrávstvuyte!} — and it comes back to you.
  \item \textbf{\ru{Спасибо}!} (\emph{Spasíbo!}) — \emph{Pozhálusta!}
\end{itemize}

With a friend:

\begin{itemize}
\raggedright
  \item \textbf{\ru{Привет}!} (\emph{Privét!})
  \item \textbf{\ru{Да}!} (\emph{Da!}) or \textbf{\ru{Нет}.} (\emph{Nyet.})
\end{itemize}

Only the greeting changes between the scenes. That one choice is the whole of the register so far.

\subsection*{Your turn}
Say these aloud:

\begin{itemize}
\raggedright
  \item the stranger's scene, both voices — \emph{Zdrávstvuyte!} … \emph{\ru{Спасибо}!} … \emph{Pozhálusta!}
  \item the friend's scene — \emph{\ru{Привет}!} — then yes, then no: \emph{\ru{Да}!} … \emph{\ru{Нет}.}
\end{itemize}

\subsection*{Script — ten letters, three kinds}
\begin{itemize}
\raggedright
  \item \textbf{False friends}: \textbf{\ru{в} \ru{р} \ru{с} \ru{н}} look like Latin \emph{B, P, C, H} and say \textbf{v, r, s, n}.
  \item \textbf{A quiet false friend}: \textbf{\ru{и}} looks like a reversed \emph{N} and says \textbf{ee}.
  \item \textbf{New shapes}: \textbf{\ru{б}} (a flag on top) says \textbf{b}; \textbf{\ru{д}} (two feet below the line) says \textbf{d}.
  \item \textbf{Honest letters}: \textbf{\ru{п} \ru{е} \ru{т}} say \textbf{p, ye, t}.
\end{itemize}

\emph{Read it:} the ten letters aloud by their sounds — \ru{в} \ru{р} \ru{с} \ru{н} \ru{и} \ru{б} \ru{д} \ru{п} \ru{е} \ru{т}

\subsection*{Writing — from sound alone}
\emph{Cover:} the letter list above and the key below, so no model and no romanization stays in view

\emph{Hear:} \emph{v, r, s, n, ee} — \emph{b, d} — \emph{p, ye, t}

\emph{Write it:} one Cyrillic letter for each sound, ten in all

\emph{Hear:} \emph{privét} — \textquotedblleft{}hi\textquotedblright{} — and \emph{nyet} — \textquotedblleft{}no\textquotedblright{}

\emph{Write it:} both words

\emph{Check:} against the key — \textbf{\ru{в} \ru{р} \ru{с} \ru{н} \ru{и}}, \textbf{\ru{б} \ru{д}}, \textbf{\ru{п} \ru{е} \ru{т}}, \textbf{\ru{привет}}, \textbf{\ru{нет}}

The usual slip is the Latin hand taking over: \emph{r} for \textbf{\ru{р}}, \emph{n} for \textbf{\ru{н}}, \emph{i} for \textbf{\ru{и}}.

\emph{Write it:} one repaired letter, then stop

\begin{tcolorbox}[breakable,colback=teal!4,colframe=teal!35!black,title={Before you move on}]
Which word answers \textbf{\ru{спасибо}}? (\emph{Pozhálusta}.) Name the four false friends by their real sounds. (v, r, s, n.) Which letter wears a flag on top, and which stands on two feet? (\textbf{\ru{б}}; \textbf{\ru{д}}.) That is the chapter: a greeting for each kind of person, the courtesy after it, and ten letters on demand.
\end{tcolorbox}
